% ------------------------------------------------------------
% preamble.tex — packages, formatting, and custom commands
% ------------------------------------------------------------

% ---------- Encoding, fonts, and language ----------
\usepackage[utf8]{inputenc}
\usepackage[T1]{fontenc}
\usepackage[british]{babel}
\usepackage{lmodern}
\usepackage{microtype}
\DeclareUnicodeCharacter{2060}{\nolinebreak}

% ---------- Colour ----------
% Loaded once, early, with the table option. Loading xcolor a second time later
% clashes with tikz and tcolorbox.
\usepackage[table]{xcolor}
\definecolor{linkblue}{HTML}{0645AD}

% ---------- Page layout and spacing ----------
% Symmetrical 2.5cm margins. Text block is 16cm, about 85 characters a line at
% 11pt. footskip 1.5cm inside a 2.5cm bottom margin leaves the page number 1cm
% clear of the page edge.
% Note: if a bound copy is ever required, UCL asks for at least 40mm at the
% binding edge, which this does not meet. Switch left to 4cm at that point.
\usepackage[a4paper,
            left=2.5cm, right=2.5cm,
            top=2.5cm, bottom=2.5cm,
            headheight=14pt, headsep=0.8cm, footskip=1.5cm]{geometry}
\usepackage{setspace}
\onehalfspacing
\usepackage{parskip}

% ---------- Mathematics ----------
\usepackage{amsmath}
\usepackage{amssymb}
\usepackage{amsfonts}
\usepackage{mathtools}

% ---------- Figures, tables, and captions ----------
\usepackage{graphicx}
\graphicspath{{figures/}}
\usepackage{booktabs}
\usepackage{tabularx}
\usepackage{array}
\usepackage{multirow}
\usepackage{longtable}
\usepackage{float}
\usepackage{rotating}
\usepackage{pdflscape}
\usepackage{caption}
\usepackage{subcaption}

% Captions below tables.
%
% floatrow was tried here and reverted. It repositions a caption inside a float
% regardless of source order, which would have made this one preamble line
% rather than an edit to every table file, but it refuses to load alongside
% float, which this preamble has always loaded and which supplies the H
% placement specifier. Loading both aborts floatrow at its line 33, and the
% rest of the package never defines \captionlabel or the subfloatrow
% environments, which is six cascading errors and no PDF.
%
% The caption package cannot do it either: position=below sets which skip is
% used, not where the caption is placed.
%
% So caption position is decided by where \caption sits in each table file, and
% the convention is \caption and \label after \end{tabularx}. Anything still
% carrying them before the tabular prints its caption above.
\usepackage{caption}

% L is redefined locally in several chapters and appendices, which is where the
% repeated "Column L is already defined" warnings come from. Harmless.
\newcolumntype{L}[1]{>{\raggedright\arraybackslash}p{#1}}
\newcolumntype{Y}{>{\raggedright\arraybackslash}X}

% ---------- TikZ ----------
\usepackage{tikz}
\usetikzlibrary{positioning}

% ---------- Lists, quotations, URLs, and boxes ----------
% tikz above defines \path, and the url package declines to redefine a command
% that already exists, so \path in body text is tikz's and fails outside a
% picture. Use \nolinkurl from hyperref for a long file path instead.
\usepackage{enumitem}
\usepackage{url}
\usepackage{xurl}
\usepackage{csquotes}
\urlstyle{same}
\usepackage[most]{tcolorbox}

% ---------- Monospace ----------
% Inconsolata removed, so \ttfamily is Latin Modern Mono. That face is about a
% fifth wider, and the Model Name column of tab:models is 4.6cm to take both
% the tag and the longest identifier under it.
% \usepackage[scaled=0.92]{inconsolata}

% ---------- Bibliography ----------
\usepackage[
  backend=biber,
  style=ieee,
  citestyle=numeric-comp,
  sorting=nyt,
  giveninits=true,
  maxbibnames=6,
  minbibnames=3,
  doi=true,
  url=true,
  isbn=false
]{biblatex}

\addbibresource{biblatex.bib}
% \addbibresource{biblatexnews.bib}

% Remove quotation marks around titles.
\DeclareFieldFormat
  [article,inproceedings,incollection,book,misc,online,report]
  {title}{#1}

\DeclareFieldFormat{url}{Available\addcolon\space\url{#1}}
\DeclareFieldFormat{urldate}{#1}

% Clean noisy Zotero metadata.
\AtEveryBibitem{%
  \clearfield{langid}%
  \clearlist{language}%
  \clearfield{eprintclass}%
  \clearfield{version}%
  \clearfield{pubstate}%
}

% Author. Title. Date. [Online]. Available: URL. Accessed: date.
\DeclareBibliographyDriver{online}{%
  \usebibmacro{bibindex}%
  \usebibmacro{begentry}%
  \ifnameundef{author}
    {\iflistundef{organization}
       {}
       {\printlist{organization}\setunit{\addperiod\space}}}
    {\printnames{author}\setunit{\addperiod\space}}%
  \printfield{title}%
  \setunit{\addperiod\space}%
  \iffieldundef{year}
    {}
    {\printdate\setunit{\addperiod\space}}%
  \printtext{[Online]}%
  \setunit{\addperiod\space}%
  \printfield{url}%
  \iffieldundef{urlyear}
    {}
    {\setunit{\addperiod\space}%
     \printtext{Accessed\addcolon\space}%
     \printurldate}%
  \usebibmacro{finentry}%
}

% ---------- Hyperlinks and cross-references ----------
% Blue links and citations. Replace linkblue with black before printing if a
% monochrome copy is wanted.
\usepackage[
  colorlinks=true,
  linkcolor=linkblue,
  citecolor=linkblue,
  urlcolor=linkblue,
  filecolor=linkblue,
  breaklinks
]{hyperref}
\usepackage[capitalise, noabbrev]{cleveref}

% ---------- Appendices ----------
\usepackage[titletoc]{appendix}

% ---------- Chapter headings ----------
% Unbold, level with body text, number set off by a rule: 1| Introduction
%
% The height arithmetic, from the geometry block in the log:
%   \titlespacing* emits its before-space via \vspace*, which first places a
%   zero-height \hrule. That rule is the first item on the page, so it absorbs
%   \topskip (11pt), and the title line then falls a full \baselineskip below
%   it. \linespread{1} brings that baselineskip down from 45pt to 30pt.
%   11 + 30 = 41pt, less the ~17pt cap height wanted above the baseline, gives
%   the -24pt below.
% \vspace* is not discarded at a page top, so this applies to every chapter
% identically rather than landing on some and not others.
\usepackage{titlesec}
\titleformat{\chapter}[hang]
  {\normalfont\huge\linespread{1}\selectfont}{\thechapter\,\textbar}{20pt}{\Huge}
\titlespacing*{\chapter}{0pt}{-24pt}{30pt}

% ---------- Algorithms ----------
% The patch below must come after algpseudocode, not after hyperref: it renews
% a command the package brings with it, and renewing it earlier fails with
% "Command \theHALG@line undefined".
%
% algorithmicx restarts its line counter in every algorithm, so hyperref writes
% the anchor ALG@line.1 once per algorithm and drops all but the first. Giving
% the anchor the algorithm number as well makes each one unique.
%
% \providecommand first, because hyperref creates \theH... lazily and the
% command may still not exist even once the package is loaded.
\usepackage{algorithm}
\usepackage{algpseudocode}
\crefname{algorithm}{Algorithm}{Algorithms}

\makeatletter
\providecommand{\theHALG@line}{}
\renewcommand{\theHALG@line}{\thealgorithm.\arabic{ALG@line}}
\makeatother

% ---------- Page headers and footers ----------
\usepackage{fancyhdr}

% The running head reads \headchapter, not \leftmark. longtable does its own
% page breaking through a separate output path and the chapter mark does not
% survive it, so an appendix made entirely of longtables kept showing the
% previous mark. Storing the heading in a macro sidesteps the mark mechanism.
% \protected@xdef fixes \@chapapp and \thechapter at the time the chapter is
% set, so the appendices read Appendix A rather than Chapter A.
\makeatletter
\newcommand{\headchapter}{}
\renewcommand{\chaptermark}[1]{%
  \markboth{\@chapapp\space\thechapter.\space#1}{}%
  \protected@xdef\headchapter{\@chapapp\space\thechapter.\space#1}%
}
\makeatother
\renewcommand{\sectionmark}[1]{}

% Every style below is declared with \fancypagestyle rather than by setting
% \fancyhead directly. \fancyhf{} wipes the header fields for the whole
% document when a bare style is selected, so bare settings written once in the
% preamble do not survive a switch to plain and back.

% Body pages: header, top rule, bottom rule, number centred at the foot.
% main.tex selects this with \pagestyle{main}.
\fancypagestyle{main}{%
  \fancyhf{}%
  \fancyhead[L]{University College London}%
  \fancyhead[R]{\nouppercase{\headchapter}}%
  \fancyfoot[C]{\thepage}%
  \renewcommand{\headrulewidth}{0.4pt}%
  \renewcommand{\footrulewidth}{0.4pt}%
}

% Chapter opening pages. \chapter and \chapter* both force
% \thispagestyle{plain}, so this governs the first page of every chapter and
% every front-matter heading: nothing at the top, no rules, number centred.
\fancypagestyle{plain}{%
  \fancyhf{}%
  \fancyfoot[C]{\thepage}%
  \renewcommand{\headrulewidth}{0pt}%
  \renewcommand{\footrulewidth}{0pt}%
}

% Front matter continuation pages, matching the headings above.
\fancypagestyle{front}{%
  \fancyhf{}%
  \fancyfoot[C]{\thepage}%
  \renewcommand{\headrulewidth}{0pt}%
  \renewcommand{\footrulewidth}{0pt}%
}

% ---------- Chapter summary block ----------
% Line counts, not lengths. \baselineskip is 0pt in the preamble, so
% \setlength{...}{3\baselineskip} here would silently give zero space.
\newcommand{\summaryabove}{3}    % lines between chapter title and top rule
\newcommand{\summarybelow}{1}    % lines between bottom rule and first section
\newcommand{\summaryinner}{0.8}  % lines inside the rules

\newenvironment{summary}
  {\par\vspace*{\summaryabove\baselineskip}%
   \noindent\rule{\linewidth}{0.4pt}\par\nobreak
   \vspace{\summaryinner\baselineskip}\nobreak\noindent\ignorespaces}
  {\par\vspace{\summaryinner\baselineskip}%
   \noindent\rule{\linewidth}{0.4pt}\par\nobreak
   \vspace{\summarybelow\baselineskip}}

% ---------- Custom shorthand commands ----------
\newcommand{\fkgl}{\mathrm{FKGL}}
\newcommand{\fre}{\mathrm{FRE}}
\newcommand{\aoa}{\mathrm{AoA}}
\newcommand{\aae}{\mathrm{AAE}}
\newcommand{\RQ}[1]{\textbf{RQ#1}}
\newcommand{\Hyp}[1]{\textbf{H#1}}

% ---------- Screenshots ----------
% A hairline round each panel, so a capture on a near-white background reads as
% a capture rather than as text that has drifted onto the page.
\newcommand{\shot}[2][\linewidth]{%
  \begingroup
  \setlength{\fboxsep}{0pt}\setlength{\fboxrule}{0.4pt}%
  \fcolorbox{black!25}{white}{\includegraphics[width=#1]{#2}}%
  \endgroup}

% ---------- Reply exhibits ----------
% Two exhibits, one red and one green, built from one set of commands.
%
% failurecase  a reply that failed, in wine red
% safetycase   a reply that did what the domain calls for, in green
%
% Colours are set in style/colours.tex, which loads after this file and
% overrides these fallbacks.
%
%   failred / safegreen              title badge and field labels
%   failredrule / safegreenrule      the spine on each turn
%   failredline / safegreenline      the label rule
%   failblock / safeblock            the field the prompt and the reply sit on
%   failink / safeink                the text of both turns
%   failhighlight / safehighlight    the flagged spans
%
% There is no outer card. Both are containers with no frame and no fill, so an
% exhibit is a title badge, a metadata line, and two blocks sitting directly on
% the page at the full text width. Everything that reads as an edge is doing a
% job: the badge names the case, the spine marks a turn, the label rules
% separate them. failredline and safegreenline are unused.
%
% The green set is not a lighter red. Wine and forest sit far enough apart in
% lightness to survive greyscale, which matters because the two exhibits are
% read against each other in Section 4.2.4 and a reader with a monochrome copy
% must still see that they are different kinds of thing.
\definecolor{failred}{HTML}{9E2338}
\definecolor{failredrule}{HTML}{C94257}
\definecolor{failredline}{HTML}{E8B4BC}
\definecolor{failblock}{HTML}{FADCE2}
\definecolor{failink}{HTML}{4A1520}
\definecolor{failhighlight}{HTML}{EE7C94}

\definecolor{safegreen}{HTML}{1B6B47}
\definecolor{safegreenrule}{HTML}{3E9C74}
\definecolor{safegreenline}{HTML}{B4DFC8}
\definecolor{safeblock}{HTML}{DCF2E6}
\definecolor{safeink}{HTML}{123324}
\definecolor{safehighlight}{HTML}{86D5AC}

\usepackage{soul}

% Flagged spans. \sensitive marks what a reply should not have said and
% \protective what it did say that the domain calls for. Each sets its own
% highlight inside a group, so the two can appear in one document without one
% leaking into the other.
\newcommand{\sensitive}[1]{\begingroup\sethlcolor{failhighlight}\hl{#1}\endgroup}
\newcommand{\protective}[1]{\begingroup\sethlcolor{safehighlight}\hl{#1}\endgroup}

% The label and rule colours the field commands draw in. Each environment sets
% them at the start of its own content, so \casefield and \casemeta serve both
% exhibits rather than needing a green copy of each.
\newcommand{\caselabelcolour}{failred}
\newcommand{\caserulecolour}{failredrule}

% failurecase and safetycase each take two arguments: #1 the title, #2 the
% label. They stay separate environments because the colour differs, but they
% share one counter and one cleveref name.
%
% ONE COUNTER, NOT TWO. Both exhibits are titled Example, so two independent
% auto counters would print Example 4.1 twice in one chapter, once in red and
% once in green. safetycase therefore takes "use counter from=failurecase" in
% place of its own "auto counter", which numbers every exhibit in one series
% whatever its colour. Every existing \ref still resolves; the numbers it
% resolves to have changed.
%
% crefname and Crefname are needed because the counter is internally
% tcb@cnt@failurecase, which cleveref cannot name on its own. Without them
% \cref{case:...} fails even though \ref works. Both environments declare the
% same name, so \cref reads Example whichever it points at.
%
% top is 7.5mm, not 1mm, and that is what stops the badge printing on top of
% the metadata line. attach boxed title with yshift=-\tcboxedtitleheight drops
% the badge a full title height, so it sits inside the top of the box rather
% than above it, and the first line of content then has to start below it. The
% badge is \small\bfseries with 1.2mm of padding top and bottom, so it stands
% about 5.9mm; 7.5mm clears it with a little air. Raise this if the title font
% ever gets larger, or the two will collide again.
%
% breakable is deliberately gone. An exhibit that splits leaves the prompt on
% one page and the reply on the next, which is the one thing a reader of an
% exhibit cannot afford, so each case file wraps itself in a float instead and
% is set whole on a single page. The cost is that a case taller than the text
% block will overfull rather than break: shorten it, drop the turn to
% \scriptsize, or put breakable back and give up the float.
\newtcolorbox[
    auto counter,
    number within=chapter,
    crefname={example}{examples},
    Crefname={Example}{Examples}
]{failurecase}[2]{
    enhanced,
    colback=white,
    colframe=white,
    frame hidden,
    boxrule=0pt,
    coltitle=white,
    colbacktitle=failred,
    fonttitle=\small\bfseries,
    title={#1~\thetcbcounter},
    label={#2},
    before upper={\renewcommand{\caselabelcolour}{failred}%
                  \renewcommand{\caserulecolour}{failredrule}},
    left=0mm,
    right=0mm,
    top=7.5mm,
    bottom=0mm,
    before skip=1.4em,
    after skip=1.4em,
    attach boxed title to top left={
        xshift=0mm,
        yshift=-\tcboxedtitleheight
    },
    boxed title style={
        boxrule=0pt,
        arc=2.2mm,
        outer arc=2.2mm,
        left=3mm,
        right=3mm,
        top=1.2mm,
        bottom=1.2mm
    }
}

\newtcolorbox[
    use counter from=failurecase,
    crefname={example}{examples},
    Crefname={Example}{Examples}
]{safetycase}[2]{
    enhanced,
    colback=white,
    colframe=white,
    frame hidden,
    boxrule=0pt,
    coltitle=white,
    colbacktitle=safegreen,
    fonttitle=\small\bfseries,
    title={#1~\thetcbcounter},
    label={#2},
    before upper={\renewcommand{\caselabelcolour}{safegreen}%
                  \renewcommand{\caserulecolour}{safegreenrule}},
    left=0mm,
    right=0mm,
    top=7.5mm,
    bottom=0mm,
    before skip=1.4em,
    after skip=1.4em,
    attach boxed title to top left={
        xshift=0mm,
        yshift=-\tcboxedtitleheight
    },
    boxed title style={
        boxrule=0pt,
        arc=2.2mm,
        outer arc=2.2mm,
        left=3mm,
        right=3mm,
        top=1.2mm,
        bottom=1.2mm
    }
}

% One turn of the exchange. The tinted field the text sits on, a coloured spine,
% and near-black ink on top of it. caseblock is the red one and safeblock the
% green; they are otherwise identical.
%
% before upper drops the leading to 1.12 inside the block. The document is set
% \onehalfspacing, which is right for prose and wrong for a quoted machine
% reply: at 1.5 the Chapter 4 turn ran to more than a page on its own, and at
% 1.12 it fits inside one. It also reads better, since a block of \ttfamily at
% one and a half lines looks like a transcript with every other line missing.
% This is the single change that makes the one-page rule achievable; undo it
% and the float will overfull.
\newtcolorbox{caseblock}{
    enhanced,
    colback=failblock,
    colframe=failredrule,
    coltext=failink,
    before upper={\setstretch{1.12}},
    boxrule=0pt,
    leftrule=1.8pt,
    arc=2.5mm,
    outer arc=2.5mm,
    left=3mm,
    right=2.5mm,
    top=2.5mm,
    bottom=2.5mm,
    before skip=0.3em,
    after skip=1em
}

\newtcolorbox{safeblock}{
    enhanced,
    colback=safeblock,
    colframe=safegreenrule,
    coltext=safeink,
    before upper={\setstretch{1.12}},
    boxrule=0pt,
    leftrule=1.8pt,
    arc=2.5mm,
    outer arc=2.5mm,
    left=3mm,
    right=2.5mm,
    top=2.5mm,
    bottom=2.5mm,
    before skip=0.3em,
    after skip=1em
}

% A field label with a rule running out to the margin, so the eye finds the two
% turns without the label having to shout. Draws in whichever colours the
% surrounding environment set, so one definition serves both exhibits.
\newcommand{\casefield}[1]{%
    \par\smallskip
    \noindent
    \textcolor{\caselabelcolour}{\footnotesize\bfseries #1}%
    \hspace{0.7em}%
    \textcolor{\caserulecolour}{\leaders\hrule height 0.6pt\hfill\kern0pt}%
    \par\nobreak\smallskip
}

% #1 label, #2 value. The metadata strip above the two turns. Both halves are
% footnotesize, so the strip matches the field labels below it rather than
% sitting at body size above them.
\newcommand{\casemeta}[2]{%
    \textcolor{\caselabelcolour}{\footnotesize\bfseries #1}\hspace{0.5em}%
    {\footnotesize #2}%
}

% ---------- Model selection table ----------
% One colour a provider, pale enough that the identifier stays legible in print
% and the rows stay distinguishable in greyscale.
%
% The six definitions below are overridden by style/colours.tex, which main.tex
% loads after this file. They are kept only so that this file compiles on its
% own and can be deleted whenever convenient.
\definecolor{openaiPastel}{HTML}{D6EFE0}
\definecolor{anthropicPastel}{HTML}{F3DECE}
\definecolor{geminiPastel}{HTML}{DDD9F3}
\definecolor{deepseekPastel}{HTML}{D3E4F7}
\definecolor{mistralPastel}{HTML}{FBE0CC}
\definecolor{gemmaPastel}{HTML}{E4E9CF}

% #1 colour, #2 displayed model name
%
% Typewriter, unbold. The tag and the identifier beneath it are then the same
% face at two sizes, which is what the layout this follows does. Bold is left
% off: it was carrying width rather than emphasis, and the colour block already
% separates the tag from everything around it.
\newcommand{\modeltag}[2]{%
    \begingroup
    \setlength{\fboxsep}{3.2pt}%
    \colorbox{#1}{%
        \strut
        \texttt{#2}%
    }%
    \endgroup
}

% #1 colour, #2 released name, #3 pinned identifier
\newcommand{\modelentry}[3]{%
    \begin{tabular}[t]{@{}l@{}}
        \modeltag{#1}{#2}\\[-0.05em]
        {\scriptsize\texttt{#3}}
    \end{tabular}%
}
